% =====================================================================
%  CH 1-4 : Vision · Couche I et taxonomie · Couches et traçabilité · Carte
% =====================================================================
\section{Vision et objectifs}\label{sec:vision}

Le développement logiciel est entré dans l’ère \emph{agent-first} : une part croissante du
code est écrite, relue et modifiée par des modèles de langage plutôt que par des humains.
Cette bascule remet en question une hypothèse centrale de nos outils — que le code source
est optimisé pour la lecture humaine. TRIADE est la réponse disciplinée à ce réexamen :
plutôt que de supprimer les langages (une impossibilité de principe, puisqu’un modèle de
langage manipule des symboles textuels et que tout format binaire rendu lisible par un
agent constitue de facto un langage), TRIADE déplace le langage vers une fine couche de
spécification maintenue par l’agent, tandis que l’artefact livré est du code machine natif
pur. Le compilateur n’est jamais livré : la notation scellée, elle, l’est — ce que
l’agent écrit est ce qui s’exécute, et rien d’autre.

Le nom encode l’architecture : trois couches indissociables — Intention (I), Binaire (B),
Carte (C) — unies par un verrou croisé cryptographique. La couche I contient le quoi et le
pourquoi : contrats typés, pré et post-conditions, invariants, cas d’échec, vecteurs de
vérification. La couche B contient le comment : des octets scellés, versionnés, à adresses
stables. La couche C est l’index naviguable qui lie les deux et sert d’unique interface de
développement à l’agent.

La v0.4.0 tire les enseignements de la première validation à l’échelle — le recodage
complet de RoomEq (démonstration v0.3) — et les convertit en normes : ce que la
démonstration a pratiqué (reverse-specs, pré-versions, partage de fixtures, mesure avant
optimisation) devient exigence ; ce que son analyse des défauts a révélé devient
taxonomie ; ce que son étude de cas établit et ne prouve pas reste explicitement hors
périmètre de preuve.

La v0.5.0 est la révision du retour du terrain. Les jalons M0 (chaîne artefact), M1
(mini-libc scellée) et M2 (patch à chaud en production, avec retour arrière) sont
franchis — chacun sur validation terrain Windows 11, chacun scellé par une balise
finale immuable (0.2.0, 0.3.0). La campagne M2 a coûté huit classes de défauts
invisibles à l’émulation (§\ref{sec:taxonomieX}) — le prix d’une leçon : un oracle
valide un artefact contre un modèle d’environnement, et la complétude du modèle ne se
démontre qu’en exécutant le binaire sur le système réel. Ce que le terrain a appris
devient norme : la boucle d’évidence (§\ref{sec:terrain}), la taxonomie X,
l’instrumentation contractualisée — et l’ouverture du jalon M3, contrats renforcés
sur entrées non fiables.

\subsection{Fonction objectif hiérarchisée — invariant de révision}\label{sec:objectif}

\begin{center}
\large\bfseries Sécurité $>$ Performance $>$ Optimisation $>$ Coût en tokens
\end{center}

\medskip

Arbitrage lexicographique : un objectif n’est jamais sacrifié au profit d’un objectif de
rang inférieur. La v0.4.0 consigna cette hiérarchie comme \textbf{invariant de révision}
(D19) ; l’édition académique la \textbf{requalifia POLITIQUE DE PROJET EXPLICITE}
(D182) — non un invariant de révision : la hiérarchie demeure la règle du projet
(l’ordre, les rangs, la non-compensation), et le document le dit dans les termes de
l’ingénierie.

\begin{table}[htbp]
\centering\small
\caption{La fonction objectif et ses métriques de validation.}
\label{tab:objectif}
\begin{tabularx}{\textwidth}{@{}l Y Y@{}}
\toprule
\textbf{Objectif} & \textbf{Signification concrète} & \textbf{Métrique de validation} \\
\midrule
Sécurité & Sûreté mémoire validée par oracle ; intégrité carte $\leftrightarrow$ binaire ;
durcissement au chargement & 0 corruption détectée ; tout échec de verrou bloque la
session ; échec au lancement = arrêt propre \\
Performance & 100 \% code machine natif ; zéro abstraction intermédiaire &
Comparaison par bloc scellé face à une référence compilée \\
Optimisation & Descente à l’ISA sur hot paths identifiés — jamais sans mesure (R11) &
Gain mesuré et journalisé par bloc optimisé \\
Coût en tokens & Édition ciblée via la carte ; jamais de lecture intégrale du binaire &
Budgets par type d’opération (§\ref{sec:carte}) \\
\bottomrule
\end{tabularx}
\end{table}

\subsection{Périmètre et non-objectifs}\label{sec:perimetre}

Cinq non-objectifs. Premièrement, TRIADE ne prétend pas remplacer les langages dans
l’industrie : c’est une architecture de laboratoire. Deuxièmement, la furtivité face aux
EDR n’est pas recherchée : l’outil cible est un utilitaire de pentest pour tests autorisés
en laboratoire, appelant les API du système via les interfaces standards du profil.
Troisièmement, l’universalité est structurelle mais progressive : chaque profil doit être
instancié et validé (M6, M7) avant d’être déclaré supporté — le protocole ne promet pas la
portabilité, il la vérifie. Quatrièmement, aucune couche d’abstraction à l’exécution
(bytecode portable, machine virtuelle, traduction à la volée) n’est introduite : elle
trahirait la fonction objectif Performance. Cinquièmement, la couche de gouvernance
validée par la démonstration et le cœur natif du protocole (notation, verrou
cryptographique sur octets, patch à chaud) restent des niveaux de preuve distincts : les
jalons M0 à M7 demeurent la charge de validation du cœur natif.

\subsection{Ce que change la v0.4.0}\label{sec:nouveautes}

Dix-sept évolutions, toutes tracées dans la correspondance du §\ref{sec:correspondance} et
dans le journal des versions, toutes rattachées aux décisions D19 à D24 : normalisation de
la couche I, taxonomie des défauts, traçabilité I $\leftrightarrow$ C $\leftrightarrow$
B $\leftrightarrow$ tests, norme de fragmentation, localisation par intention, R10
renforcée, R11 nouvelle, suite de conformité par profil, parité inter-profils,
pré-versions, reverse-specs, journal d’optimisation, modèle de menace étendu, journal des
décisions normalisé, tableau de bord de conformité. Aucune règle existante n’est
affaiblie ; aucune couche n’est ajoutée à l’exécution ; la hiérarchie de la fonction
objectif est inchangée.

\subsection{Ce que change la v0.5.0}

Huit évolutions, toutes tracées dans la correspondance (§\ref{sec:evolutions050}) et le
journal des versions, toutes rattachées aux décisions D26 à D45 : la validation
terrain devient une couche de vérification normée — boucle d’évidence, journal
d’exécution, instrumentation contractualisée, diff chirurgical (§\ref{sec:terrain}) ;
la taxonomie des défauts s’étend des classes T (contrat non écrit) aux classes X —
défauts invisibles à l’émulation, huit classes nommées par la campagne M2
(§\ref{sec:taxonomieX}) ; l’état des jalons est constaté — M0, M1 et M2 franchis sur
validations terrain Windows 11, balises finales 0.2.0 et 0.3.0 posées
(§\ref{sec:jalons}) ; le budget de localisation est corrigé à 64 tokens (D31) et les
budgets affichent leurs valeurs mesurées ; les volumes d’oracle effectifs sont
consignés ; les preuves terrain sont archivées brutes avec leurs empreintes ; le
corpus de démonstration (huit itérations instrumentées) devient l’occurrence de
référence de la méthode ; le jalon M3 est ouvert — contrats renforcés sur entrées non
fiables et fuzzing à propriétés (D45). Aucune règle n’est modifiée ; aucune couche
n’est ajoutée à l’exécution ; la hiérarchie de la fonction objectif est inchangée.

\subsection{Ce que change la v0.5.1}\label{sec:nouveautes051}

Un composant nouveau et deux constats, tous tracés dans la correspondance
(§\ref{sec:evolutions051}) et le journal des versions, rattachés à la décision D75 : le
\textbf{coffre-fort des phrases maîtresses} — conteneur chiffré TRIADE-CF1 centralisant
toutes les phrases maîtresses par projet, sur la pile cryptographique de la carte au
repos (Argon2id aux planchers, HKDF à contexte dédié, XChaCha20-Poly1305), doté d’un
autotest de falsification et de règles d’export/import (continuité,
§\ref{sec:gouvernance} ; FORMAT\_CARTE § 11) ; l’état des jalons est constaté — M3
franchi le 28 octobre 2026 sur la campagne million (D72), la conformité permanente
(D73) et le scellage de la balise finale 0.4.0 (D74) ; le modèle de menace gagne la
ligne du vol du coffre-fort (§\ref{sec:menaces}). Aucune règle n’est modifiée ; aucune
couche n’est ajoutée à l’exécution ; la hiérarchie de la fonction objectif est inchangée.

\subsection{Ce que change la v0.5.2}\label{sec:nouveautes052}

Une évolution de discipline, tracée dans la correspondance (§\ref{sec:evolutions052}) et
le journal des versions, rattachée à la décision D81 puis \textbf{amendée par D223} :
la phrase de contrôle fut mise sous garde durable (le fichier au dépôt, du 4 au
19 novembre 2026) ; depuis le séquestre vécu D183, la garde ACTIVE est les CINQ
PARTS Shamir à seuil trois — la reconstitution à froid quand la phrase sert,
la phrase en AUCUN dépôt (public depuis D220) ; tout geste qui tenterait de
réécrire la phrase (ou son fichier) au dépôt est refusé avec son nom. Le constat fondateur : la
rotation D80 avait produit une phrase neuve dont la remise prévue en fin de tour n’a
jamais eu lieu — la session d’exécution est morte avant (pannes d’infrastructure
récurrentes) — et la phrase ne vivait plus que dans le bac à sable éphémère : un décès
d’infra de plus et le coffre entier devenait indéchiffrable à jamais. L’arbitrage est
consigné : le risque de \textsc{perte} (vécu, répété) est jugé supérieur au risque
d’\textsc{exposition} (URL non devinable) ; le modèle de menace en tire la conséquence
honnête — la confidentialité des phrases au repos repose désormais sur le contrôle
d’accès au dossier du coffre Drive, la mitigation première devient la restriction du
partage par lien, et une phrase entrée dans l’historique git y demeure à vie : on ne
compte plus sur sa confidentialité, la rotation reste le remède. L’état des jalons est
constaté — M4 ouvert le 28 octobre 2026 (D77, couche I écrite-avant-code : huit
spécifications \code{tls} et deux révisions \code{http}). Aucune règle n’est modifiée ;
aucune couche n’est ajoutée à l’exécution ; la hiérarchie de la fonction objectif est
inchangée.

\subsection{Ce que change la v0.6.0}\label{sec:nouveautes060}

La révision de génération attendue par le plan : le \textbf{programme de souveraineté
cryptographique est accompli et prouvé sur terrain} — décision de passation D133,
arbitrage du développeur. Trois axes livrés (§\ref{sec:souverainete}) : l'\textbf{axe 2}
(D86) — le coffre CF2 à clé maîtresse répartie par Shamir, cinq gardes seuil trois,
sentinelle à double seuil puis PORTE D120 ; l'\textbf{axe 3} (D87) — le sceau de
l'Intention : dix arbres de couche I signés par la clé Ed25519 hors-ligne du
développeur, gardes de build et de patch ; l'\textbf{axe 4} (D121, né du terrain) — le
coffre CF3 : la serrure à triple scellement, le piège sur média séparé, la sentinelle
instantanée \emph{par invocation} et le cycle en deux actes. L'\textbf{axe 1} (D85,
attestation TEE) demeure une exigence consignée : l'infrastructure sans racine de
confiance matérielle ne peut pas la satisfaire et le dit. La \textbf{chaîne réelle du
coffre} a migré sur le terrain du développeur les 30--31 octobre 2026 — CF1 vérifié
INTÈGRE, phrase de contrôle changée (D110), CF2 migré (cinq gardes, seuil 3, trois
copies hors-ligne), CF3 migré (époque 1, checksum \code{0cc89a96}), verdict à froid
VERT deux fois — les cérémonies étant conduites par le maître de cérémonie graphique
(D107--D118) avec le kit Drive tenu identique au miroir par construction (D130). Un
\textbf{coffre neuf naissait CF3} : \code{init} direct à l'époque 1 — la chaîne
de migration est l'histoire du coffre réel, pas un rite obligé (la v0.6.1 remet ce
choix au développeur — §\ref{sec:nouveautes061}, §\ref{sec:initniveau}). Un document
complémentaire, le \emph{Livre du Coffre} (même date, même dépôt), porte la posture
opératoire du maître : cérémonies pas à pas, runbook, mode dégradé. Aucune règle du
cœur n'est modifiée ; aucune couche n'est ajoutée à l'exécution ; la hiérarchie de la
fonction objectif est inchangée.

\subsection{Ce que change la v0.6.1}\label{sec:nouveautes061}

Une seule règle, explicitement invalidée et remplacée (§\ref{sec:evolutions061}) :
le niveau de sécurité d'un coffre neuf. La v0.6.0 (D133) l'avait fixé — un coffre
neuf naît CF3, \code{init} direct à l'époque 1 ; la v0.6.1 (D137) en fait
\textbf{l'arbitrage du développeur au moment de créer son projet de développement} :
CF1, CF2 ou CF3, chacun avec son acte de naissance, ses obligations et son guide
(§\ref{sec:initniveau}) — imposer CF3 par défaut risque d'être trop contraignant
selon le projet : média du piège à produire dès le premier jour, cycle en deux actes
à chaque session, sachet physique à tenir — la contrainte se justifie quand le secret
la mérite. La chaîne de migration CF1 $\to$ CF2 $\to$ CF3 demeure disponible à tout
moment pour un coffre dont les besoins grandissent ; le coffre réel de TRIADE reste
CF3 (époque 1, checksum \code{0cc89a96}). Aucune autre règle n'est modifiée ; aucune
couche n'est ajoutée à l'exécution ; la hiérarchie de la fonction objectif est
inchangée.

\subsection{Ce que change la v0.7.0}\label{sec:nouveautes070}

La révision de consolidation du programme d'après M7 (D160) : ce que la clôture de la
v0.6.1 tenait « hors périmètre courant » entre dans la spécification, sur résultats
mesurés — la clause est levée NOMMÉMENT pour les deux premières familles, les profils
additionnels demeurent hors périmètre. Le \textbf{jalon M8} (le bloc optionnel de
syscalls directs — le domaine \code{sys}) est franchi : la charte à huit règles (le SSN
est une donnée, jamais un octet de code ; le motif du stub vérifié avant toute
extraction ; un stub divergent est un refus nommé — le résolveur ne devine jamais, il ne
s'adapte jamais en silence), les spécifications sys-01…03 écrites-avant-code, le harnais
différentiel 61 contrôles verts sur fixtures \code{ntdll} synthétiques (le laboratoire
ne connaît jamais le SSN du terrain), la cérémonie de l'ère passée avant le build (12
arbres, D162/D163), le run terrain sur le VRAI \code{ntdll} — la twin de la ligne
hello-02 par les DEUX chemins (l'IAT \code{kernel32} et le syscall direct
\code{NtWriteFile}), le refus $-4$ nommé du stub crocheté — recoupé 30/30 sur pièces
croisées, la finale \code{0.9.0} scellée (Entrée 164, D165 : le scellage exige le
verdict du recoup comme pièce ; le scellage est un fait de carte, l'artefact ne bouge
jamais). Le \textbf{jalon M9} (TLS 1.3 pur en blocs scellés — le domaine \code{tls13},
l'ère 1.0.0) est ouvert (D166) : la couche I complète — 28 spécifications en six
familles (sha, hkdf, x25519, aead, enregistrement, poignée, principal), l'instrument
\code{verifier\_specs\_m9a.py} 117 contrôles verts (D167) — et la cérémonie de l'ère
passée le 05/11 : la remise du développeur, TREIZE arbres signés hors ligne à la clé du
04/11 (la continuité prouvée par le déterminisme Ed25519, D163 : les onze arbres
inchangés rendent les signatures de l'ère M8), le garde repassé VERT 13/13 — la dette
D166 soldée sur pièces. Le \textbf{sceau de l'Intention} porte désormais treize arbres
(§\ref{sec:sceauI}) et les cérémonies d'ère entrent au protocole
(§\ref{sec:ceremoniesere}) : la dette se nomme à la naissance, jamais après coup (D166,
règle 8) ; la charge se prépare au geste qui sert la main du développeur (D168). Aucune
règle du cœur n'est modifiée ; aucune couche n'est ajoutée à l'exécution ; la hiérarchie
de la fonction objectif est inchangée.

\subsection{Ce que change la v0.9.0}\label{sec:nouveautes090}

La révision de consolidation du programme d'après M10 (D178/D179). Le \textbf{jalon
M10} (le test grandeur nature — le domaine \code{loupe}) est \textbf{franchi} :
l'application « Loupe », l'éditeur hexadécimal au panneau d'empreintes, est NÉE
ENTIÈREMENT DE LA MÉTHODE — la cérémonie d'ère passée sur pièces (le garde à
QUATORZE arbres, la dette D166/D171 soldée), la couche I de 18 spécifications
écrites-avant-code (D172), les QUATRE chantiers de notation (D173–D175 : les 18
blocs NOTÉS-MESURÉS-PROUVÉS, 159 contrôles verts aux trois graines), le kit déployé
au coffre \code{/test} (M10-(d)) — ET DOUZE RUNS TERRAIN (M10-(e), Entrées 202–222) :
ONZE DÉFAUTS NOMMÉS SUR ONZE RUNS FAUTIFS (le douzième vert de bout en bout —
le résidu [T0-PÉRIMÉ] divulgué, non compté : une bannière de texte, aucune
porte touchée), ONZE RETOUCHES DOCUMENTÉES — les refus du chargeur
0xc000007b (l'enveloppe prouvée de la lignée NORME DE FAIT), [CADRE-NÉ]
(l'initialiseur \code{.data} MIROIR EXACT du \code{CADRE\_VIERGE} du harnais — LE
GARDE GUÉRI : le contrôle dérive du plan, jamais de la pièce qu'il contrôle), la
classe des arguments pile (l'audit d'ABI msx64 NÉ AU BUILD), [SESSION] $-7$
VERROUILLÉ (L'OMBRE au kit, l'artefact acquitté), [HINSTANCE-PEB] (l'entrée
réécrite : r12 = GetModuleHandleW(NULL) LE VRAI MODULE), [D54-QUATRIÈME] +
[EAX-SANS-SIGNE] (l'IAT relogée à 0x3100 en QUATRE plages séparées par NULL, les 20
rendus négatifs étendus au 64 bits), [OFT-QWORD] (les QUATRE OFT au pas DWORD du
chargeur — la sixième mort NOMMÉE par le journal d'événements [4/4]),
[FREE-SANS-TYPE] (le HMODULE en POINTEUR, l'essai isolé pour TOUTE classe
d'exception), [RET-NU] (\code{mov eax, 1 ; ret} — le verdict ATTACH VRAI PAR
CONSTRUCTION), [SONDE-EXE] (le faux négatif STRUCTUREL de la sonde de liaison : le
chargeur de bibliothèque ne résout JAMAIS les imports d'un module exécutable — LA
LIAISON PROUVÉE PAR LE VIVANT) — ET LA VITRINE A VÉCU (0 et 0 aux dixième et
onzième runs : L'APPLICATION NÉE DE TRIADE A VÉCU CHEZ SON UTILISATEUR).

LE ONZIÈME RUN recoupé 39/39 : \textbf{LA BALISE \code{loupe 0.1.0} SCELLÉE PAR
FAIT DE CARTE} (D178 : la dette de la carte soldée par la RÉGÉNÉRATION DES
NOTATIONS (nov. 2026) — 18/18 notations == la pièce OCTET POUR
OCTET, l'écart de l'audit
$587 \to 0$ ; la cristallisation $26 \times 10/10$ ; la re-signature à la graine
durable D144 — LE TÉMOIN NE BOUGE PAS ; l'anti-rollback ARMÉ ; \code{loupe.exe}
INTACT \code{054de960…} — LE MÊME binaire que le run a couru) : \textbf{LA MÉTHODE
TRIADE ENTIÈRE PROUVÉE PAR LA CRÉATION — de l'intention (la couche I) à l'octet
vécu chez son utilisateur}.

\textbf{LE PROGRAMME D'APRÈS M10 APPARTIENT AU DÉVELOPPEUR} : la version
\code{0.1.1-dev.1} est COMMANDÉE (D179) — la dette [R10-PÉRIMÉE] (les 17 sommes de
l'ère dev.5 que les retouches des Entrées 210/214 changèrent sans rafraîchir la
table) SOLDÉE À LA NAISSANCE : la pièce naît de LA RECOMPOSITION même (chaque
somme == son bloc relu, par construction), le qword version porté 0x00010101,
0 OCTET AILLEURS — le code BYTE-IDENTIQUE au vécu des douze runs (R4 : le contrat
stable ; R7 : les adresses intactes), la carte GÉNÉRATION 17 ; le kit
\code{0.1.1-dev.1} DÉPOSÉ au coffre \code{/test} (loupe-kit-13, le réceptacle
ARMÉ) — LE DOUZIÈME RUN VÉCU ET RECOUPÉ 39/39 (l'Entrée 222 : T0$\to$T7 tout
vert pour la cinquième fois, la cascade 205 vécue pour la deuxième, la vitrine
vivante la troisième — 0 et 0, aucune porte touchée) : LA BALISE
\code{loupe 0.1.1} SCELLÉE PAR FAIT DE CARTE sur run terrain recoupé (D181 :
la dette [R10-PÉRIMÉE] soldée à la naissance, la cristallisation 26 × 10/10,
l'anti-rollback armé — \code{loupe.exe} INTACT \code{c1115105…}, LE MÊME
binaire que le douzième run a couru). Le projet fils COFFRE-FORT demeure dégelable au rythme du
développeur (D135 amendée par D171) : la frontière fondatrice demeure INTACTE.
Aucune règle du cœur n'est modifiée ; aucune couche n'est ajoutée à l'exécution ;
la hiérarchie de la fonction objectif est inchangée.

\subsection{Ce que change la v0.10.0}\label{sec:nouveautes100}

La révision de L'ÉDITION ACADÉMIQUE (D182, la commande du développeur du
11/11 : trois fichiers d'instructions au dossier \code{/instructions} du
coffre). La présente révision \textbf{consigne} l'édition académique du
protocole ; elle ne modifie AUCUNE règle, n'ajoute AUCUNE couche
d'exécution, laisse la hiérarchie de la fonction objectif inchangée — le
modèle D180 : le protocole consigne, il ne légifère pas.

\textbf{LE MANUSCRIT ACADÉMIQUE} (\code{docs/TRIADE\_manuscrit\_v0.10.0.pdf},
la chaîne \code{scripts/manuscrit\_v100/}) : une édition IMRAD en anglais
(résumé français inclus), structureée autour de QUATRE questions de
recherche — le taux de défauts échappés (onze classes nommées sur douze
runs terrain), le coût en tokens par changement validé (51 / 719 / 1 088
mesurés M1, D31/D34), les classes de défauts invisibles à l'émulation (les
cinq familles : chargeur/image, ABI/SSE, sémantique du chargeur de
bibliothèque, conventions de codes d'erreur, hypothèses du harnais), la
sûreté du patch à chaud (zéro régression post-patch, monotonie signée,
anti-rollback armé) ; l'auteur humain unique (l'agent IA décrit en
méthodes et remerciements, JAMAIS co-auteur) ; les déclarations
Ethics/dual-use, Data Availability, Limitations (N=1, circularité de
l'oracle, absence de baseline, sécurité mémoire non prouvée, temps
constant non analysé) ; chaque donnée du manuscrit est GÉNÉRÉE depuis le
journal par \code{scripts/extraire\_donnees\_v100.py} — le manuscrit
n'affirme rien qui ne soit extrait et vérifié.

\textbf{LA FORMULATION CORRIGÉE} (le manuscrit et, par lui, le protocole) :
l’ancienne formule sémantique (le compilateur réputé retranché du
pipeline) devient « AUCUNE chaîne de traduction
optimisante ni opaque ; assemblage WYSIWYG vérifié par round-trip » ;
« code volé inexploitable » devient « l'évolution du code volé est
bloquée, son exécution ne l'est pas » ; les majuscules emphatiques et le
vocabulaire sacralisant cèdent aux termes standards (key ceremony, signed
manifest, sealed build) ; l'invariant D19 (Sécurité > Performance >
Optimisation > Coût en tokens) est requalifié POLITIQUE DE PROJET
EXPLICITE — non un invariant de révision : la hiérarchie demeure la règle
du projet, et le document le dit dans les termes de l'ingénierie.

\textbf{LA PROPOSITION D183 (documentée, NON appliquée au coffre réel)} :
le séquestre chiffré de la phrase de contrôle — l'alternative à la garde
durable D81 (la phrase en clair au dépôt et au coffre Drive, l'arbitrage
PERTE > EXPOSITION) : la phrase chiffrée vers une clé hors-ligne, ou
partagée en parts Shamir à seuil, la phrase en clair RETIRÉE du dépôt ;
l'analyse comparative chiffrée des risques perte/exposition vit dans la
décision D183 ; le coffre réel du développeur EST PASSÉ au séquestre
à seuil le 19 novembre 2026 (D220/D223 : la cérémonie vécue, la
phrase répartie, le dépôt passé public — la matrice de visibilité
D219 honorée) — l'énoncé d'époque (« demeure à la garde D81 ») est
l'histoire datée du document, jamais l'état courant.

\textbf{L'ANNEXE A PORÉE À SHA-256} (D184) : les empreintes de gouvernance
de l'annexe A (les modules du coffre) passent de md5 à sha-256, calculées
au contenu canonique LF du dépôt (la même méthode, l'empreinte longue) ;
md5 demeure UNIQUEMENT une fonctionnalité d'affichage de Loupe (le panneau
d'empreintes, loupe-06) — AUCUN octet de la pièce ne change, et les
empreintes md5 des preuves historiques (le journal, les archives) demeurent
ce qu'elles sont : les preuves ne se réécrivent jamais.

\textbf{LES LIVRABLES} : \code{TRIADE\_soumission\_v0.10.0.zip} (cinq
fichiers au plus, readme inclus, le manuscrit en pdf — prêt pour un comité
de révision) et \code{TRIADE\_complet\_v0.10.0.zip} (l'intégralité :
sources LaTeX, bibliographie, tableaux de données, script d'extraction,
extraits de journal, empreintes). La v0.9.0 demeure : ses sources et son
PDF archivés à \code{docs/}, la chaîne \code{spec\_v090} intacte — chaque
édition vit à sa chaîne.

\subsection{Ce que change la v0.11.0}\label{sec:nouveautes110}

La révision des PREUVES RENFORCÉES (D185/D186, la commande du développeur :
le deuxième fichier d'instructions). AUCUNE règle modifiée, AUCUNE couche
d'exécution ajoutée, hiérarchie inchangée ; les instruments nouveaux
VIVENT AU LABORATOIRE uniquement (l'exigence de la commande), aucun
artefact scellé ne bouge, aucun arbre signé ne bouge.

\textbf{LES SIX INSTRUMENTS} (les graines consignées, reproductibles —
\code{scripts/preuves\_v110/}) : (1) LE MUTATION TESTING — 832 mutants
d'un octet des blocs scellés (copies de laboratoire), tués par les
instruments EXISTANTS dans l'ordre du coût : 694/832 (83,4 \%), les
familles cryptographiques à 100 \%, les 138 survivants CONCENTRÉS dans
les blocs d'orchestration NOMMÉS (la limite mesurée, jamais cachée) ;
(2) LE FUZZING GUIDÉ PAR COUVERTURE — loupe-01 (le chemin non fiable,
l'OS virtuel, les canaris) : 60,5 \% de couverture du bloc, zéro crash,
zéro violation de canaris ; http-10 (la pièce 0.4.0 scellée) : 358
instructions couvertes, zéro faute nommée ; (3) LA SÛRETÉ MÉMOIRE
STATIQUE BORNÉE — l'extension de l'audit msx64 : 2~098 accès classés,
1~181 vérifiés au cadre (zéro hors-cadre : la classe du stack-smash
ABSENTE), 463 non résolus NOMMÉS, les limites consignées (analyse
linéaire, aucun alias) ; (4) LE TEMPS CONSTANT — le dudect
D'ÉMULATION : toutes les familles crypto CONSTANTES (traces
indépendantes du secret — x25519 : 1~238~959 instructions pour TOUTE
clé), le temps réel du terrain NON MESURÉ (la limite explicite) ;
(5) LES BENCHMARKS COMPARATIFS — la référence C à contrat équivalent
(gcc 14.2 -O2/-Os, la référence clang ABSENTE du labo : l'écart nommé),
LE MÊME émulateur pour les deux camps : sha256 4~524 instructions/op
(blocs) contre 69~132 (gcc -O2) — un facteur 15,3 — pour 765 o contre
5~436 ; md5/aes ÉCARTÉS par l'AUTO-VÉRIFICATION de l'instrument (le C
de référence fautif REFUSÉ : jamais un chiffre faux) ; (6)
L'INDÉPENDANCE DE L'ORACLE — les voies distinctes (dérivation décimale,
os.urandom, spec ↔ harnais ↔ hashlib) : ZÉRO divergence.

\textbf{LA SPÉCIFICATION CSPRNG (D186)} : la couche I manquante de la
source d'aléa COMBLÉE AU LABORATOIRE — le brouillon
\code{tls13-29-source-csprng.md} (le générateur = le CSPRNG de l'OS
hôte via le kit, le reseed = la politique du noyau, l'échec
FAIL-CLOSED, le KIT96 seule interface) ; l'entrée à l'arbre signé
EXIGE la cérémonie d'ère du développeur (le quinzième arbre, le modèle
D177) — jusqu'alors le brouillon est la référence de laboratoire.

\textbf{LES VECTEURS RFC 8446} (\code{tls13/fixtures/rfc8446.py}) : les
constantes normatives natives (les sentinelles DOWNGRD, la valeur
magique HRR, les types d'enregistrement, la suite et le groupe) posées
au corpus de fixtures partagées avec leur paragraphe d'origine — le
plan RECOISÉ (SUITE, GROUPE, HRR : conformes).

\textbf{LE MANUSCRIT v0.11.0} : la section Results étendue (le § 3.5 du
manuscrit — les six instruments, la table générée), les limites MISSES à
JOUR (la sûreté mémoire BORNÉE pas prouvée, le temps constant
STRUCTUREL pas réel, la réplication par second agent à venir), la
conclusion réécrite sur CE QUI RESTE OUVERT précisément mesuré.

\subsection{Ce que change la v0.12.0}\label{sec:nouveautes120}

La révision du MODÈLE FORMEL ET DE LA GOUVERNANCE ASSAINIE (D187/D188/D189,
le troisième fichier d'instruction). AUCUNE règle modifiée, AUCUNE couche
d'exécution ajoutée, hiérarchie inchangée.

\textbf{LE MODÈLE FORMEL DE CF3} (D187) : la chaîne du coffre scellé
modélisée en machine à états explicite et vérifiée par EXPLORATION
EXHAUSTIVE (154 états, 1~152 transitions, époques ≤ 3, adversaire à seuul
≤ 2 éléments avec rejeu d'époques passées) — LES QUATRE PROPRIÉTÉS
PROUVÉES DANS LE MODÈLE : P1 le secret sous seuil (0 contre-exemple),
P2 le refus du roll-back (la monotonie structurelle, constatée
mécaniquement transition par transition), P3 l'absence de destruction sur
phrase erronée (0 contre-exemple), P4 le rayon borné à une époque (0
violation) ; LES HYPOTHÈSES EXPLICITES A1--A5 (le modèle symbolique, le
seuil, le média, les anciens couples, le HORS MODÈLE : canaux
auxiliaires, implémentation, trois éléments, quantique) ; LE
VÉRIFICATEUR NOMMÉ : Tamarin/ProVerif ABSENTS du laboratoire —
l'explorateur d'états explicites tient la place, le script EST le modèle
(reproductible d'un clone frais).

\textbf{LE MODÈLE DE MENACE EN DEUX PLANS} (D188) : le tableau unique
restructuré selon la distinction fondatrice du protocole — le PLAN IP
(la propriété intellectuelle : rétro-ingénierie, expropriation,
redistribution — la confidentialité d'abord, les pertes REMÉDIABLES) et
le PLAN SÛRETÉ D'EXÉCUTION (la corruption, le rollback, l'environnement
— l'intégrité et la disponibilité d'abord, les pertes IRRÉMÉDIABLES) ;
L'ORDRE INTRA-SÉCURITÉ EXPLICITE (la politique de projet) : quand les
deux plans se rencontrent, le plan sûreté GAGNE (un coffre cassé ne
protège rien) ; LA TENSION D81 consignée avec la contre-proposition D183
et l'analyse chiffrée (l'espérance de perte annuelle ≈ 4,0 ×
10\textsuperscript{−2}·C\textsubscript{exposition} sous D81 contre ≈ 1,0
× 10\textsuperscript{−3} sous le séquestre Shamir — un facteur ≈ 40) —
D81 DEMEURE la décision active, le geste appartient au développeur.

\textbf{LA CRISTALLISATION ÉVALUÉE} (D189, CORRIGÉE par D190) : les
caractéristiques opérationnelles de la règle réelle (sceller ssi $n$
rebuilds consécutifs identiques) : P(sceller$|q$) = $(1-q)^n$ — à
$q = 50$\,\% : 3,1 \% ($n{=}5$), 0,098 \% ($n{=}10$), soit 32× plus
strict (les quantiles unilatéraux exacts 1,86σ [valeur exacte
1,863]/3,10σ) ; à $q = 1$\,\% :
la règle scelle un build sain avec probabilité 95,1 \%/90,4 \% — les
valeurs « α = 3,29·10\textsuperscript{−2} / 1,08·10\textsuperscript{−3} »
de D189 étaient des RAPPORTS DE VRAISEMBLANCE $(0,5/0,99)^n$ mal
étiquetés (la correction D190) ; LE SEUIL DEMEURE 10 (la politique de projet, jamais changée
en silence) — l'analyse est la JUSTIFICATION consignée et la voie
documentée d'un relâchement futur ; LE PROTOCOLE DE RÉPLICATION proposé
(les trois axes : second développeur, second domaine, second profil ;
la grille de codage à cinq classes terrain + trois classes laboratoire
— étiquetée T/X à l'époque, renommée F1--F5/L1--L3 par D203 pour
lever la collision avec les taxonomies de défauts T1--T6/X1--X8 du
§\ref{sec:taxonomie} ; la
double annotation ; le κ de Cohen ≥ 0,61 ; les critères d'acceptation
pré-enregistrés) — PROPOSÉ, jamais exécuté : l'arbitrage appartient au
développeur.

\textbf{LES MANIFESTES DE KIT À SHA-256} (D187) : le manifeste du kit de
cérémonie (maj\_kit.py) consigne désormais sha256 — la suite de la
migration D184 ; les kits DÉJÀ LIVRÉS demeurent (leurs manifestes md5
sont des PREUVES, les preuves ne se réécrivent jamais) ; md5 demeure à
l'idempotence Drive (l'API donne le md5) et à l'affichage Loupe
(fonctionnalité).

\textbf{LE MANUSCRIT v0.12.0} : la section § 4.1 (le modèle formel et
ses quatre propriétés avec leurs statuts, le modèle de menace en deux
plans avec l'ordre intra-sécurité, la cristallisation, le protocole de
réplication), la table des preuves étendue (v0.10.0 + v0.11.0 + v0.12.0).

\subsection{Ce que change la v0.13.0}\label{sec:nouveautes130}

LA RÉVISION DES CORRECTIONS DE L'ÉDITION (D190/D191/D192) — la commande du
développeur : une liste exhaustive d'incohérences vérifiées une par une, des
statistiques reformulées, la neutralité documentaire, un benchmark audité, un
état de l'art étendu. AUCUNE règle du cœur modifiée, AUCUNE couche d'exécution
ajoutée, AUCUN artefact scellé touché — l'édition CORRIGE les documents, elle
ne légifère pas (le modèle D180/D182).

\textbf{LES INCOHÉRENCES CORRIGÉES} (la liste de la commande, vérifiée une à
une) : les quatorze arbres signés (la cohérence avec le garde 14/14) ; la
référence au manuscrit accompagnant portée v0.13.0 ; le compte de lignes du
journal UNIFIÉ en valeur exacte extraite (les macros générées par le script
d'extraction — jamais un arrondi) ; « Fig.~5 » devient « Table 5 » ; la Data
Availability portée aux archives v0.13.0 ; la Table 6 des preuves RÉGÉNÉRÉE
intégralement en anglais par le générateur ANTI-TRONCATURE (les cellules ne
sont jamais coupées, chaque ligne est PARSÉE depuis les sorties brutes des
instruments) ; la note TOTAL/brouillons au tableau des budgets ($8\,774 +
5\,099 = 13\,873$, les brouillons exclus) ; le paragraphe de consolidation
du § 1.10 ALIGNÉ sur D181 (douze runs, la balise \code{loupe 0.1.1} scellée) ;
le tableau de menaces : « le code volé est inexploitable » devient
« l'évolution du code volé est bloquée, son exécution ne l'est pas » ; les
renvois au manuscrit recalés (§ 3.5 les preuves renforcées, § 4.1 le modèle
formel).

\textbf{LES STATISTIQUES REFORMULÉES} (D190) : la cristallisation — les
caractéristiques opérationnelles P(sceller$|q$) = $(1-q)^n$ remplacent les
« α » de D189 (des rapports de vraisemblance mal étiquetés) ; les conversions
sigma $\sqrt{-2\ln p}$ remplacées par les quantiles normaux unilatéraux
exacts (1,86σ/3,10σ) ; l'IC95~\% EXACT de 11/12 (Clopper-Pearson
$[0,615 ; 0,998]$, avec la réserve de non-indépendance des itérations
séquentielles) ; la règle de comptage opérationnelle des défauts échappés
DÉFINIE en méthodes du manuscrit (les runs 6/7 à deux défauts comptés une
fois chacun ; le résidu [T0-PÉRIMÉ] du run 12 divulgué, non compté — et
pourquoi).

\textbf{LA NEUTRALITÉ DOCUMENTAIRE} (D190) : l'anti-debug léger par mesure
de temps DOCUMENTÉ HONNÊTEMENT dans les deux documents (la portée : la
dissuasion, jamais la résistance — le résiduel « assumé » du tableau des
menaces demeure) ; la conséquence de D81 AJOUTÉE au modèle de menace de l'époque —
DATÉE du 4 au 19 novembre 2026 (la phrase de contrôle alors au dépôt +
S\_serrure en en-tête clair = deux éléments trivialement disponibles ;
depuis D223 la phrase vit sous séquestre, le voleur de la machine tient
UN élément, la confidentialité repose sur le seuil) ; ; la frontière D105/D81 pour le clone du laboratoire (les tests
du labo tournent sur des clones temporaires à phrases de test visiblement
marquées — le conteneur réel jamais ouvert par un test) ; LE KIT DE LA
CÉRÉMONIE D183 PRÉPARÉ SANS L'APPLIQUER (\code{scripts/kit\_ceremonie\_d183/} :
l'outil au partage CF2 D120 TEL QUEL, l'autotest VERT sur phrase de test
publique — les 10 combinaisons de 3 parts exactes, les 10 paires sous le
seuil sans fuite —, le LISEZMOI pas à pas, le manifeste et les listes de
contrôle ; le geste appartient au développeur).

\textbf{LE BENCHMARK AUDITÉ} (D191) : le facteur « 15,3× » de v0.11.0
RETIRÉ — un ARTEFACT DE PÉRIMÈTRE (l'instrument lisait le compteur de la
FINALISATION SEULE après le retour du pipeline, comparé à l'opération
COMPLÈTE de la référence ; 4,16 instructions/ronde était le signe — sous
tout minimum scalaire plausible) ; les chiffres corrigés : l'opération
complète des blocs 75\,983 instructions contre 69\,132 (gcc~14.2 -O2) et
77\,505 (-Os), dans le même émulateur sur la même entrée, le condensé
vérifié AVANT le comptage et les 17 invocations de compression confirmées
mécaniquement — la notation WYSIWYG est COMPÉTITIVE (elle exécute
$+9{,}9$~\% d'instructions \emph{de plus} que gcc -O2 [75\,983 contre
69\,132] et $2{,}0$~\% \emph{de moins} que -Os [77\,505]), plus compacte que -O2 (1\,021 contre 1\,482 octets, tables
comprises), moins que -Os (548) ; AUCUNE extension ISA au-delà du baseline
des deux côtés (gcc emploie des mouvements SSE2 du baseline que la notation
GPR pure n'emploie pas — l'écart rapporté sans ajustement) ; le terrain M7
REQUALIFIÉ (qemu-riscv64 TCG au noyau réel, la réserve instrumentale D149
nommée — la sémantique matérielle I-cache de \code{fence.i} NON prouvée sur
silicium, jamais revendiquée).

\textbf{L'ÉTAT DE L'ART ÉTENDU} (D192) : LLM Compiler (Meta 2024), CompCert
(Leroy 2009), Kitsune (2012) et POLUS (2007) pour la mise à jour dynamique,
SWE-agent (2024), la lignée sim-to-real (Jakobi 1995, Peng 2018) et le
hardware-in-the-loop (Bacic 2005) — chaque référence vérifiée DOI/arXiv ;
la revendication « first quantified account » SCOPÉE (le niveau notation
scellée sous gouvernance à intention signée) ; RQ3 RECADRÉE : la
quantification des GARANTIES SILENCIEUSES DU COMPILATEUR (ce que le
compilateur fournit gratuitement et que la notation à la main doit
re-dériver, leçon nommée par leçon).

\subsection{Ce que change la v0.14.0}\label{sec:nouveautes140}

LA RÉVISION DU PROGRAMME EMPIRIQUE (D193/D194/D195, le cinquième fichier
d'instructions) : le bras de référence compilateur PRÉ-ENREGISTRÉ, la
re-mesure RQ2, l'indépendance par instance isolée, le re-jeu des
survivants de mutation. AUCUNE règle du cœur modifiée, AUCUNE couche
d'exécution ajoutée, AUCUN artefact scellé touché.

\textbf{LE PRÉ-ENREGISTREMENT} (D193) : le plan de comparaison (les
hypothèses H1-H4, les critères CA1-CA5, la méthode, les six blocs, les
vecteurs) consigné AVANT toute mesure, son empreinte sha256 verrouillée —
toute modification postérieure du plan serait une fraude attrapable.

\textbf{LE BRAS DE RÉFÉRENCE} (D195) : SIX blocs de Loupe réimplémentés
EN C depuis les MÊMES CONTRATS (les cinq familles couvertes), compilés
gcc 14.2 -O2/-Os (clang absent, l'écart nommé), les primitives OS prises
par table de pointeurs aux MÊMES shims du mini-OS, AUTO-VÉRIFIÉS contre
les vecteurs dorés AVANT toute mesure (l'autotest a attrapé deux fautes
de table K au geste — la méthode appliquée à elle-même). LA MESURE :
ZÉRO DIVERGENCE aux vecteurs (la notation scellée et la référence
compilée sont observationnellement équivalentes aux contrats) ; les
instructions à travail égal : sha-256 PARITAIRE (9~089 contre 9~146 —
l'héritage PAR LA CARTE), la lecture paritaire (48/47), les blocs purs
plus chers à la notation (641/477, 193/69, 39/16), la cascade saine
102~176/75~560 (l'adaptation du panneau nommée : le camp C calcule la
moitié sha seulement) ; les tailles MIXTES ; la mutation aux mêmes
tirages : la notation testable au MÊME DEGRÉ (H4 tenue), le camp compilé
plus fragile aux mutants d'appels OS (16/16 par le tueur de séquence).

\textbf{LA RE-MESURE RQ2} (D195) : les budgets de bout en bout sur la
carte de l'ère Loupe, distributions complètes (n=10/opération,
cl100k\_base) : localiser 580 [496..706], éditer 2~143 [1~085..3~463],
ajouter 3~510 [3~128..4~091], la baseline naive 300~304 tokens (le
désassemblage complet) — les valeurs M1 (51/719/1~088) HISTORIQUES,
jamais remplacées ; LA LECTURE DÉFAVORABLE AU PROTOCOLE consignée telle
quelle (les opérations coûtent 3-11$\times$ leur M1 — les ratios
recalculés depuis \code{remesure\_rq2.csv} : éditer $2\,143/719 =
3{,}0\times$, ajouter $3\,510/1\,088 = 3{,}2\times$, localiser
$580/51 = 11{,}4\times$ — les notations plus
longues, les contrats plus denses) ; le raisonnement ESTIMÉ
($\times$1,08, le rapport génération/contenu des retouches réelles —
l'estimation nommée, jamais une mesure).

\textbf{L'INDÉPENDANCE} (D195) : l'oracle régénéré par une instance
isolée (un contexte neuf, la dérivation depuis les définitions
mathématiques) — ZÉRO DIVERGENCE sur les 128 constantes ; L'ÉCART NOMMÉ :
l'infrastructure d'agents du labo n'expose QU'UNE famille — l'indépendance
inter-familles DEMEURE ouverte ; la voie alternative (la racine cubique
entière exacte) confirme indépendamment.

\textbf{LE RE-JEU DES SURVIVANTS} (D195) : les 832 mutants de v0.11.0
re-jeués aux MÊMES tirages avec les tueurs de COMPOSITION (la boucle
winmain complète pour loupe, la composition de la poignée +
l'enregistrement pour l'orchestration tls13, le harnais autonome pour
sys) : 776/832 tués (93,3 \% — contre 83,4 \%), tls13 et sys à
100 \%, les 56 survivants loupe NOMMÉS (les chemins hors scénario).

\textbf{LE KIT TERRAIN} (\code{scripts/kit\_terrain\_bras\_c/} : les
binaires -O2/-Os + la source + le LISEZMOI) : EN ATTENTE DU RUN TERRAIN —
les runs appartiennent au développeur, le labo ne fabrique JAMAIS un
résultat.

\subsection{Ce que change la v0.15.0}\label{sec:nouveautes150}

LES PREUVES NON BORNÉES ET LE RAFFINEMENT (le sixième fichier
d'instructions, ic09ydeli sous son vrai nom — le croisement
d'attribution d'hier RECROUPÉ et CORRIGÉ, D197) : le modèle formel
généralisé, l'équivalence modèle--implémentation vérifiée, les statuts
de preuve systématiques. AUCUNE règle du cœur modifiée, AUCUNE couche
d'exécution ajoutée, AUCUN artefact scellé touché.

\textbf{LA DÉVIATION CONSTATÉE} : Tamarin ET ProVerif demeurent ABSENTS
du laboratoire (constat re-vérifié au geste à l'ouverture de l'édition)
— la route consignée à D196 est suivie : la preuve par INVARIANT
INDUCTIF dans l'explorateur.

\textbf{L'INVARIANT INDUCTIF} (\code{scripts/preuves\_v150/
modele\_invariant\_v150.py}) : le plafond d'époque de l'explorateur
v0.12.0 (un artefact d'exploration, jamais du design — l'implémentation
fait $e \leftarrow e+1$ sans borne) est RETIRÉ ; l'invariant $I$
(I$_1$ secret $\Rightarrow$ seuil $=3$ ; I$_2$ aucune destruction
injuste ; I$_3$ média du piège et fragments CF2 intacts ; I$_4$ époque
$\ge 1$) est initialisé à l'état initial, PRÉSERVÉ par chacune des 13
règles sur l'ensemble FINI des formes d'état (288 formes, les valeurs
contraintes par $I$ énumérées dans le domaine $I$, les deux classes
d'époque $\{1\}$ et $\{\ge 2\}$), l'uniformité de la classe $\ge 2$
témoignée mécaniquement par translations $\delta \in \{7, 50, 1000\}$
aux verdicts IDENTIQUES ; la monotonie de l'époque vérifiée
INCONDITIONNELLEMENT (4\,096 formes, zéro décroissance) ; le rejeu
retourne l'état inchangé. LES QUATRE PROPRIÉTÉS P1--P4 SONT PROUVÉES
NON BORNÉES (toute époque, toute profondeur, sous A1--A4 ; A5 hors
modèle nommé). La FIDÉLITÉ du portage : l'explorateur aux bornes et au
plafond historiques reproduit EXACTEMENT D187 (154 états, 1~152
transitions) ; les bornes étendues (époques $\le 5$ profondeur 8 :
228/1~902 ; époques $\le 7$ profondeur 12 : 372/3~342) trouvent ZÉRO
état atteignable hors de $I$ — la preuve bornée CONFIRME désormais la
preuve non bornée au lieu d'y suppléer.

\textbf{LE RAFFINEMENT} (\code{scripts/preuves\_v150/
raffinement\_v150.py}) : chaque règle abstraite REJOUÉE contre
l'implémentation Python (un coffre de TEST à phrase publique marquée,
D192 — le coffre réel n'est jamais ouvert par un test) : le cycle
honnête sur CINQ époques VÉCUES aux deux niveaux (la borne v0.12.0
dépassée), la phrase erronée aux deux actes (quinze essais, zéro
destruction), la falsification du conteneur (une destruction LOCALE,
un procès-verbal, le média intact), la falsification du piège (refus
nommé, rien détruit), le rejeu d'époque passée en TROIS variantes (le
miroir local, les époques croisées, le témoin hors-bande), le piège
absent (IMPOSSIBLE, pas refusé), le piège d'un autre coffre, et
l'adversaire sous le seuil (32 triplets de moitiés forgées, zéro
passage du verif) : \textbf{34 pas de raffinement, ZÉRO ÉCART} — les
chemins d'échec inclus ; trois gestes hors modèle consignés comme
tels.

\textbf{LES TROIS KITS DE MESURE MATÉRIELLE} : le dudect temps réel
(quatre cibles extraites de la pièce scellée, l'appel natif VALIDÉ au
vecteur FIPS « abc », le mesureur C double backend dont l'autotest
POSIX n'éprouve que le mécanisme), le \code{fence.i} sur silicium (le
scénario M7, la cible changée), le terrain tiers Loupe (T0$\to$T7+S0
intégral) — TOUS \emph{EN ATTENTE DE RUN MATÉRIEL} : aucun chiffre
n'existe, aucun n'est fabriqué.

\textbf{LA RÉPLICATION} : la grille F/L FINALISÉE (les cinq classes
terrain F1--F5 et trois laboratoire L1--L3, définitions
opérationnelles, exemples vécus — l'ancienne grille T/X renommée par
D203, la table de correspondance au kit), le calculateur $\kappa$ de Cohen AUTOTESTÉ (trois cas
synthétiques au $\kappa$ calculable à la main), les critères
pré-enregistrés D189 INCHANGÉS ; le second annotateur À RECRUTER, le
second domaine EN ATTENTE d'une vraie commande — les statuts consignés.

\textbf{LES STATUTS SYSTÉMATIQUES} (la commande, point 4) : chaque
affiliation de preuve porte son statut — \emph{prouvé non borné} /
\emph{prouvé borné} / \emph{en attente de run matériel} — la table
générée depuis les pièces (\code{scripts/statuts\_preuves\_v150.py}),
AUCUNE revendication matérielle tant que le run n'est pas recoupé.

\textbf{LA DETTE DOCUMENTAIRE GUÉRIE} : la table des preuves du
manuscrit (29 lignes en v0.13.0/v0.14.0) n'était JAMAIS RENDUE — le
flot dépassait la page et \LaTeX{} le rejetait silencieusement (« Float
too large », le constat de l'édition) : la table devient LONGTABLE,
les 34 lignes sont rendues, toutes les éditions sont relues.

\textbf{LA GOUVERNANCE} (D197) : le croisement d'attribution des
instructions n°5/n°6 recoupé aux métadonnées Drive (les deux fichiers
déposés 11/11 13:21 UTC, jamais modifiés — ysddxpk0o porte le programme
v0.14.0 exécuté, ic09ydeli porte le programme v0.15.0 ouvert) et
corrigé aux gestes (ysddxpk0o archivé, ic09ydeli rappelé d'ARCHIVE
vers \code{/instructions} — l'instruction active sous son vrai nom) ;
LA RÈGLE DU CONTENU CITÉ (à l'ouverture de toute instruction : lecture
intégrale, empreinte confrontée, le contenu cité à l'entrée qui ouvre
l'exécution) ; LA RÈGLE DU DÉVELOPPEUR (12/11) : à chaque livraison
d'une nouvelle version d'un document, les anciennes versions sont
archivées (racine $\to$ \code{/ARCHIVE}, déplacement vérifié — la
racine du coffre ne porte que les versions courantes).

\subsection{Ce que change la v0.15.1}\label{sec:nouveautes151}

UNE ÉDITION DE COHÉRENCE INTER-PIÈCES (D218, la première instruction
de la vague du 18/11 — 376f15f18) : AUCUNE règle modifiée, aucune
couche ajoutée, aucun artefact scellé touché. Cinq réconciliations sur
pièces, chacune sourcée : (a) le §\ref{sec:jalons} porte le douzième
run VÉCU ET SCELLÉ (D181, l'Entrée 222 — la balise \code{loupe 0.1.1}
scellée sur le run vert recoupé 39/39, l'état attendu est clos) ;
(b) le §\ref{sec:objectif} dit la requalification D19 (D182) :
POLITIQUE DE PROJET EXPLICITE, non invariant de révision ; (c) le
§\ref{sec:sceauI} compte les QUATORZE arbres signés (D177 — l'arbre
\code{loupe} entré à la table canonique) ; (d) l'annexe décisionnelle
porte son titre réel D75--D198 (la table allait déjà à D198) ; (e) les
budgets du §\ref{sec:carte} sont annotés HISTORIQUES (ère de
conception, v0.3/D31) contre les mesures d'édition v0.17--v0.18
(D195). La décision D218 (append-only) consigne les corrections et
leur sourcement ; le linter inter-pièces confronte les valeurs
partagées aux cinq pièces du pack à chaque build (D218).

\subsection{Ce que change la v0.15.2}\label{sec:nouveautes152}

L'ÉDITION DE LA TRANSITION PRÉPARÉE (D219, l'instruction jx4l5vy5c — la
vague du 18/11, la deuxième) : AUCUNE règle modifiée, aucune couche
ajoutée, aucun artefact scellé touché, le coffre réel et les phrases
réelles JAMAIS approchés (D105 : passeur.exe/passeur-c.dll re-téléchargés
intacts, le sceau 14 arbres verts). La ligne D81 du modèle de menace
(§\ref{sec:menaces}) est portée à SON STATUT EXACT : la transition vers
le séquestre chiffré D183 est \textbf{préparée, geste au développeur} —
le kit de cérémonie rafraîchi (l'autotest rejoué : 10/10 combinaisons
exactes, 10/10 paires sous le seuil inertes, la falsification PROUVÉE —
une part forgée REFUSÉE par la porte de l'empreinte de contrôle ; le
mode \code{--verifier} fail-closed né au geste), l'analyse chiffrée
perte/exposition reprise (le facteur $\approx$ 40 en faveur du
séquestre), le scan de secrets de l'historique git complet au dry-run
vert (l'équivalent Python nommé, gitleaks/trufflehog absents — six
secrets plantés sur six attrapés, zéro faux négatif ; le vrai dépôt
trié : la phrase D81 connue, neuf fixtures de test nommées, aucune
fuite réelle au-delà de l'arbitrage assumé), la checklist de rotation de
toute phrase exposée, la matrice de décision de visibilité du dépôt
(privé tant que la cérémonie n'est pas vécue et la rotation vérifiée ;
public + Zenodo seulement après), et l'extension transparence publique
DOCUMENTÉE NON APPLIQUÉE (l'ancrage du témoin d'époque à Rekor/Sigstore
ou RFC 3161 : le format du témoin, la vérification machine, le delta au
modèle de menace aux deux bords). D81 DEMEURE LA DÉCISION ACTIVE
jusqu'à ce que la cérémonie soit vécue de bout en bout — le protocole
prépare, le maître décide (D87).

\subsection{Ce que change la v0.15.9}\label{sec:nouveautes159}

La v0.15.9 est la révision de la \textbf{publication de bascule et de la racine de confiance} (D253,
D254). Elle corrige une promesse fausse : aucun trampoline de douze à seize octets ne se publie par une
« écriture alignée unique atomique » — l’atomicité d’un magasin naturel s’arrête à huit octets, et
seize octets n’existent que par un unique \texttt{LOCK CMPXCHG16B} sous CPUID (§\ref{sec:primitives}).
Elle sépare ensuite l’intégrité, l’authenticité et la fraîcheur de l’artefact, nomme l’ancre réelle par
profil et ses limites, et borne ce que R10 garantit : le rejeu du seul binaire est refusé par le plancher,
la restauration cohérente d’un couple binaire + état échappe à toute ancre locale
(§\ref{sec:racine}). Aucune valeur mesurée n’est modifiée.

\section{La couche I : format normalisé des spécifications}\label{sec:coucheI}

La v0.3 décrivait la couche I comme « contrats légers » (D3), sans en fixer le gabarit —
la légèreté demeure, la structure devient normative (D19). Une spécification I n’est pas un
document libre : c’est un artefact diffable dont la structure conditionne la vérification.

\subsection{Identifiants stables}

Chaque spécification reçoit à sa création un identifiant \code{<DOM>-<NN>} : un domaine
parmi une liste figée par projet et un numéro séquentiel. L’identifiant est immuable et
jamais réutilisé ; un module supprimé laisse son identifiant vacant à jamais. C’est
l’équivalent exact, côté intention, de l’adresse stable côté binaire (R7) : les références
croisées — dépendances, carte, tests — ne cassent jamais. Le couple (identifiant,
révision) est la clé de traçabilité du §\ref{sec:traceabilite}.

\subsection{Gabarit obligatoire}

\begin{lstlisting}
<DOM>-<NN> . <nom court> . revision <R>
Intention : <une phrase — le pourquoi>
Contrat :
  pre  : P1..Pn    (preconditions numerotees, une phrase chacune)
  post : Q1..Qn    (postconditions numerotees)
  inv  : I1..In    (invariants)
Echecs :
  E1 : <condition> -> <sanction nommee : exception | arret propre | repli>
Vecteurs :
  V1 : <vecteur dore | propriete differentielle | test de rejet>
Fixtures    : <references aux fixtures contractuelles partagees>
Provenance  : ecrite-avant-code | reverse-spec de <binaire scelle @ tag>
Decisions   : <Dxx...>
\end{lstlisting}

Sept blocs, tous obligatoires — les échecs et vecteurs ne peuvent être vides que pour un
module sans effet observable, situation rare à justifier dans le champ Intention. Trois
règles de rédaction conditionnent la diffabilité exigée par R2 : \textbf{une exigence par
phrase numérotée} (le diff montre exactement quelle exigence change) ; \textbf{aucune
connaissance dupliquée} (renvoi par identifiant, sinon défaut T5) ; \textbf{chaque
pre/post/invariant relie au moins un vecteur} (une exigence invérifiable est une exigence
mal formulée). La provenance distingue les spécifications écrites avant le code — le cas
normal, exigé par R2 — des reverse-specs (§\ref{sec:reversespecs}).

\subsection{Révision et cycle de vie}

La révision s’incrémente à chaque changement du contrat. Une évolution qui rompt un
contrat existant ne modifie pas la spécification — elle en crée une nouvelle, l’ancienne
étant dépréciée avec renvoi : miroir exact du cycle de vie des blocs. Les appelants
dépendent du couple (identifiant, révision) ; la carte porte ce couple dans chaque entrée.

\subsection{Fixtures contractuelles partagées}

Quand plusieurs consommateurs se lient au même contrat externe — sortie d’un outil
système, format de fichier, API distante — les échantillons qui définissent ce contrat
appartiennent à la spécification, pas aux consommateurs. Le champ Fixtures référence un
jeu unique ; toute copie privée est un défaut de classe T5. Le mécanisme est né dans la
démonstration RoomEq : l’installateur scellé et l’application recodée partagent les
fixtures de parsing de pnputil — le défaut historique qui a survécu à trois versions
vivait précisément de la duplication.

\subsection{Taxonomie des défauts de contrat non écrit}\label{sec:taxonomie}

La classe de défauts que le protocole existe pour éliminer a un nom : \emph{contrat non
écrit}. La v0.4.0 en fixe la taxonomie (D19), née de l’analyse des défauts historiques de
la démonstration.

\begin{table}[htbp]
\centering\small
\caption{Taxonomie des défauts de contrat non écrit.}
\label{tab:taxonomie}
\begin{tabularx}{\textwidth}{@{}l Y Y@{}}
\toprule
\textbf{Classe} & \textbf{Définition} & \textbf{Verrou} \\
\midrule
T1 — valeur-seuil enfouie & Constante critique sans nom ni justification & constante
nommée + rationale dans la spec I \\
T2 — convention d’interface implicite & Structure ou ABI déclarée de mémoire & contrat
typé + invariant de taille testé \\
T3 — dépendance à l’hôte & Comportement dépendant de l’OS d’exécution & contrat
indépendant de l’hôte, union fixée \\
T4 — parsing de surface fragile & Analyse de libellés localisés ou volatils & fragments
indépendants des libellés + fixtures partagées \\
T5 — contrat dupliqué & Même connaissance codée en N exemplaires divergents & source
unique : spec I + fixtures \\
T6 — garantie non testée & Comportement critique sans vecteur de rejet & cas d’échec
nommé + test de rejet \\
\bottomrule
\end{tabularx}
\end{table}

Deux règles d’usage : tout défaut corrigé est classé dans cette taxonomie au journal des
décisions (la métrique « défauts ouverts par classe » du tableau de bord en dérive) ; la
taxonomie sert de liste de contrôle de rédaction — chaque bloc du gabarit verrouille au
moins une classe. Les occurrences de référence viennent de la démonstration : le seuil de
détection non justifié (T1+T6), la structure \code{SP\_DEVINFO\_DATA} de 424 octets au
lieu de 32 (T2, échec \code{0x6F8}), la sanitisation de noms dépendante de l’hôte (T3), le
parsing de pnputil sur libellés français (T4+T5, trois copies).

\section{Les trois couches : architecture et traçabilité}\label{sec:couches}

Le diff d’intention descend de I vers C, le verrou croisé relie C et B dans les deux sens,
et la couche B seule quitte le laboratoire. \textbf{Aucune couche n’est optionnelle} —
c’est une propriété normative : privé de I, le système perd son ancrage sémantique ; privé
de C, l’agent devient aveugle ; privé de B, il n’y a plus rien à exécuter. Toute
implémentation qui livrerait B sans I et C à jour viole R1 ; toute évolution qui
contournerait C pour éditer B directement viole R3 et la boucle.

\begin{table}[htbp]
\centering\small
\caption{Les trois couches.}
\label{tab:couches}
\begin{tabularx}{\textwidth}{@{}l l Y Y@{}}
\toprule
\textbf{Couche} & \textbf{Visibilité} & \textbf{Contenu} & \textbf{Rôle} \\
\midrule
I & Privée & Spécifications normalisées : contrats pré/post, invariants, cas d’échec,
vecteurs & Compression de l’intention \\
C & Privée, signée, chiffrée & Index des blocs, adresses, preuves, manifestes,
traçabilité & Unique interface de développement \\
B & Livrée & Code machine scellé, padding, checksums runtime & Exécution pure,
sémantiquement muette \\
\bottomrule
\end{tabularx}
\end{table}

La séparation croise une seconde dimension : la cible. La couche I est indépendante de la
cible — un contrat s’écrit une fois, ses types abstraits étant paramétrés par le profil.
La couche B est par profil. La couche C indexe l’ensemble : chaque entrée porte son
profil, et frag-bld organise les empreintes par profil. Quatre invariants structurels,
vérifiés à chaque cycle — leur violation est un refus de build : B n’est jamais nu (R1) ;
les adresses exposées sont stables (R7, par profil) ; la signature couvre la carte entière
(R9) ; B est auto-vérifiant au chargement (R10 renforcée).

\subsection{Traçabilité I $\leftrightarrow$ C $\leftrightarrow$ B $\leftrightarrow$
tests}\label{sec:traceabilite}

Chaque entrée de la carte référence quatre maillons : la spécification (identifiant +
révision), l’entrée d’index et le contrat embarqué, le(s) bloc(s) instancié(s) avec
empreinte(s) par profil, et les vecteurs exécutés avec le verdict du dernier cycle. Trois
propriétés sont contrôlées mécaniquement par la \textbf{commande d’audit de traçabilité} :
aucune spécification sans entrée de carte (l’exception visible étant la reverse-spec en
statut « prête ») ; aucun bloc scellé sans spécification référencée ; aucun vecteur
orphelin. Une chaîne incomplète est un refus de build — au même titre qu’un binaire nu.
La traçabilité est la condition de l’audit par des tiers : elle permet de répondre à
« quelle exigence ce bloc implémente-t-il, et quel test la défend ? » sans lecture du
binaire.

\section{La carte : fragmentation et budgets}\label{sec:carte}

\begin{sloppypar}
La carte est le cœur opérationnel du protocole : elle rend le binaire navigable, porte les
preuves et concentre la valeur intellectuelle. Elle est autodescriptive : un en-tête
normalisé consigne la version du protocole, la suite cryptographique~— hachage, signature,
chiffrement, chacun dans une liste blanche (D15)~— et les profils actifs. Sa spécification
détaillée fait l’objet de \code{FORMAT\_CARTE.md} — fragments, format d’entrée, verrou,
cycle de vie, export, re-cryptage —, mise à jour avec la présente révision : l’en-tête
porte \code{version=0.4.0} et chaque entrée de bloc porte le champ de traçabilité
\code{spec=<DOM>-<NN>:R}.
\end{sloppypar}

\paragraph{Norme de fragmentation.} La carte est fragmentée en quatre — frag-loc, frag-con,
frag-prf, frag-bld — et la v0.4.0 pose la norme : aucune opération de la boucle n’exige la
lecture de plus d’un fragment ; \textbf{aucune opération n’exige la lecture intégrale du
binaire} — le balayage est interdit (§\ref{sec:boucle}) ; l’audit complet de la carte est
une opération exceptionnelle, non budgétée, journalisée.

\paragraph{Budgets par type d’opération.}

\begin{table}[htbp]
\centering\small
\caption{Budgets en tokens par opération — valeurs HISTORIQUES de l'ère de conception (localiser amendé à 64 tokens par D31 après la mesure M1 — le gabarit canonique v0.4.0 mesurait 43 tokens en cl100k\_base ; autres budgets inchangés depuis la v0.3) ; les MESURES d'édition v0.17--v0.18 (l'instrumenté au tokenizer natif :
localiser 1~277, éditer 3~374, ajouter 3~674 — puis le pointeur seul adopté
OPT-04/D222 : localiser 1~103 mesuré à la graine consignée ; les valeurs
580/2~143/3~510 sont la re-mesure v0.14.0 DE D195, au contenu seul, jamais
confondues avec l'instrumenté) font foi pour l'état courant.}
\label{tab:budgets}
\begin{tabularx}{\textwidth}{@{}Y Y l@{}}
\toprule
\textbf{Opération} & \textbf{Données lues} & \textbf{Budget} \\
\midrule
Localiser un bloc & frag-loc — une entrée & 64 tokens max (D31) ; la réponse AU POINTEUR SEUL depuis la v0.15.4 (D222 : l'ID + l'empreinte, le client expande) \\
Éditer simplement un bloc & frag-loc + frag-con + notation & 3 000 tokens max \\
Ajouter un bloc & frag-loc + frag-con + seuils frag-prf & 6 000 tokens max \\
Instancier un profil nouveau & I partagée + frag-con + notation du profil & 6 000 tokens
max par bloc \\
Audit complet (rare) & carte entière déchiffrée & non budgété \\
\bottomrule
\end{tabularx}
\end{table}

Indicateurs de santé : précision de localisation (cible 95 \% dès M0), coût moyen d’une
édition simple (cible 5 000 tokens), taux de régression post-patch (cible zéro).
Instrumentés depuis M1 : la mesure effective (51 / 719 / 1 088 tokens, D31) a corrigé
le budget de localisation à 64 et valide les deux autres ; le tableau de bord les
affiche comme métriques mesurées, avec leur écart.
\subsection{Ce que change la v0.15.3}\label{sec:nouveautes153}

L'ALIGNEMENT DU MODÈLE SUR LA RÈGLE (l'instruction 5ey7gmjra — la vague
du 18/11, la troisième) : AUCUNE règle modifiée, aucune couche ajoutée,
aucun artefact scellé touché. La règle de la sentinelle CF3 disait ce
qu'elle dit depuis la v0.12.0 : \textbf{seule la falsification
\emph{prouvée} détruit} — destruction locale, jamais le média ni les
fragments. Le MODÈLE FORMEL v0.15.0 sous-approximait cette règle : la
garde de \code{falsifie\_conteneur} exigeait l'état canonique, la
campagne générative (10~000 traces, la graine consignée) croisa l'écart
660 fois — la falsification pendant une session ouverte : le modèle
refusait, la sentinelle concrète détruisait (la découverte consignée,
réduite par ddmin à la séquence minimale [ouvrir, falsifier]). La
v0.15.3 \textbf{ferme l'écart} : la garde devient le piège intègre
seul — la falsification prouvée offerte aux DEUX phases, l'effet
inchangé (destruction locale, PV, média et fragments intacts) — et le
programme de preuve est re-déchargé sur le modèle étendu : les
invariants inductifs (Init $\Rightarrow$ I, I $\wedge$ Next
$\Rightarrow$ I$'$) à entiers non bornés NoError, l'alignement
TLA+~$\leftrightarrow$ sémantique Python 13/13 règles (garde, image,
activation — SMT non borné), le lemme de translation d'époque re-prouvé
(LT1/LT2, contrôle négatif attrapé), l'espace atteignable borné
INVARIANT (la chaîne de fidélité TLC : 154 états distincts aux bornes
historiques — seules les transitions tirées croissent, 896 $\to$ 926),
et la campagne re-jouée \textbf{à ZÉRO divergence} : 30~000 traces aux
trois graines (la consignée 20261112 + deux nouvelles), chaque chemin
d'échec couvert, le rejeu bit-identique deux fois — contre le registre
v0.15.0 : exactement les 660 traces de la classe d'écart changées, rien
d'autre. L'instrument temps constant est promu au \textbf{hash des
traces complètes} (une seule empreinte par cible sur 120 mesures — le
compte unique demeure, subsumé), et le kit de mesure matérielle
PRÉ-ENREGISTRÉ est livré en attente du geste développeur (n$\ge$10
exécutions par N, deux machines au moins, la fréquence fixée, le pic
tls13-02/tls13-13 à N=200 pré-enregistré pour investigation) — aucun
chiffre matériel fabriqué.

\subsection{Ce que change la v0.15.4}\label{sec:nouveautes154}

LE POINTEUR SEUL (l'instruction cq01qv1za — la vague du 18/11, la
quatrième) : UNE RÈGLE ADDITIVE à l'opération \emph{localiser}, aucune
règle existante modifiée, aucun artefact scellé touché. Le format de
réponse de localiser devient le \textbf{pointeur seul} : l'agent rend
UNE ligne — l'identifiant de l'entrée de carte et l'empreinte
\code{sha256} de la notation (deux champs LITTÉRAUX du matériau,
jamais un calcul — la leçon OPT-02 : le champ calculé est le champ
faussé). Le \textbf{client expande} le pointeur auprès de la carte :
recherche de l'identifiant, confrontation de l'empreinte octet pour
octet, expansion de l'entête complet, confrontation au bloc soumis —
QUATRE refus nommés fail-closed (POINTEUR-INCONNU,
EMPREINTE-DIVERGENTE, EXPANSION-DIVERGENTE, HORS-FORMAT). La mesure
OPT-04 (la même graine que OPT-02, 18 appels réels au tokenizer natif,
le cl100k parallèle, la re-tokenisation indépendante à ±1 token) : le
total 1~714 → 1~103 natifs (−35,7~\%), la réponse 770 → 52
(−93,3~\%), (nov. 2026) zéro divergence 18/18, l'expansion client
verte 18/18. Le
budget D31 (64 tokens à la première ligne) demeure : le pointeur le
tient avec marge. La décision D222 (append-only) consigne l'adoption.



\subsection{Ce que change la v0.15.5}\label{sec:nouveautes155}

L'édition v0.15.5 est L'ÉDITION DE L'ÉTAT VÉCU DE LA GOUVERNANCE
(l'instruction ansz8lvdf, la décision D223) : AUCUNE règle R1--R11
modifiée, AUCUN profil touché — le modèle de menace (§ 15) et la
gouvernance (§ 14) portés à l'État VÉCU : (1) la cérémonie de séquestre
D183 EST VÉCUE (19 novembre 2026) — le découpage Shamir cinq parts
seuil trois, la répartition sur cinq médias dont deux hors nuage, la
reconstitution à froid PROUVÉE DEUX FOIS, la phrase de contrôle commise
TOURNÉE et l'ancienne REFUSÉE à la ré-ouverture ; (2) la ligne D81 du
modèle de menace porte le séquestre chiffré comme LA GARDE ACTIVE — la
phrase de contrôle ne vit plus au dépôt ni au coffre Drive ; (3) le
dépôt est PUBLIC (D220) — la matrice de visibilité honorée ; (4) la
rotation VÉRIFIÉE : le re-scan de l'historique git complet re-essayé et
trié (dix constats, aucune fuite réelle au-delà de l'arbitrage assumé).
Les corrections de cohérence de l'instruction gq117wphb portées : le
titre de l'annexe décisionnelle à la plage réelle de la table
(D75--D223), l'ordre croissant des sections de versions (§ 1.18--§ 1.21),
l'étiquette D195 (la re-mesure v0.14.0 au contenu seul) DISTINGUÉE des
mesures instrumentées v0.17 et de la médiane adoptée OPT-04. Le Livre
du Coffre v1.4 porte le runbook § 7.5 au statut VÉCU ; le manuscrit
v0.19.6 porte les statuts à l'état vécu ; le linter inter-pièces v196
confronte le statut D81/D183/D220 aux cinq pièces.
